\subsubsection{避孕套的正确使用、尺寸选择与储存}

前面"使用方法"给出了标准步骤，但临床与真实世界里的失败（破裂、滑脱）大多来自几个可避免的细节。本节把最容易出错的环节集中讲清楚。

\vspace{0.5em}
\noindent\textbf{1. 尺寸：不是"均码就好"}

\begin{itemize}
	\item 避孕套按标称宽度（阔度）分号，常见为 49--56 mm。\textbf{过紧}易在射精时被撑破或勒痛，\textbf{过松}则容易滑脱，两者都会显著提高失败率。
	\item 选择标准：佩戴后应贴合但不留勒痕，能顺利卷至阴茎根部，勃起后无明显束缚感。若反复出现"用后滑脱"或"中途破裂"，优先怀疑尺寸不符，而非"运气差"。
	\item 市售"加厚""超薄""带凸点"等属体验差异，与安全性无必然关系；\textbf{真正决定安全的是正确佩戴与全程使用}。
\end{itemize}

\noindent\textbf{2. 佩戴：四个高频错误}

\begin{itemize}
	\item \textbf{忘了捏储精囊}：不排出前端空气，射精时压力会把套子撑破。正确做法是\textbf{用指腹捏住前端小囊}再向下展开。
	\item \textbf{戴反后再翻过来用}：若展开时发现方向反了，必须\textbf{换一只新的}——翻面使用时套外已可能沾上前列腺液（含精子），等同未戴。
	\item \textbf{用牙或剪刀开包装}：极易划破。用手沿锯齿边撕开即可。
	\item \textbf{用油性润滑剂}：凡士林、婴儿油、按摩油、食用油会\textbf{在数十秒内溶解乳胶}，造成破裂。需润滑时应选\textbf{水基或硅基}产品；若用聚氨酯套，则不受油基影响。
\end{itemize}

\noindent\textbf{3. "全程使用"而非"中途戴"}

\begin{itemize}
	\item 射精前的尿道球腺液（俗称"前列腺液"）中\textbf{可能已含少量精子}，且可携带多种病原体。因此"先做一段再戴上"既不能可靠避孕，也无法预防性传播感染。
	\item 正确做法是\textbf{从任何插入性接触开始前就佩戴}，直至射精后、阴茎仍勃起时\textbf{按住套口根部}一并抽出，避免精液外溢。
	\item 一次只能戴一只；\textbf{双层套}因乳胶相互摩擦反而更易破裂，不推荐。
\end{itemize}

\noindent\textbf{4. 储存与有效期}

\begin{itemize}
	\item 乳胶套怕热、怕光、怕油。应存放于\textbf{阴凉干燥处}，避免长期放在钱包、车内手套箱或贴身口袋（体温与摩擦会加速老化）。
	\item 使用前\textbf{检查有效期}与包装是否破损；过期乳胶套弹性下降，破裂风险上升。
	\item 若包装已干瘪、漏气，或套子发脆、发黏，直接丢弃。
\end{itemize}

\noindent\textbf{5. 万一失败：补救要做对}

\begin{itemize}
	\item \textbf{滑脱或破裂后}：尽快（数小时内）清洗外阴，但\textbf{不要冲洗阴道}；评估是否需要\textbf{紧急避孕}——左炔诺孕酮在 72 小时内、乌利司他在 120 小时内越早越有效，含铜宫内节育器在 5 天内避孕效果最佳（详见"紧急避孕"一节）。
	\item 有性传播感染暴露风险者，应在窗口期后检测，必要时咨询医生是否进行暴露后预防（如 HIV 的 PEP 需在 72 小时内启动）。
	\item \textbf{不要靠"事后冲洗"或"立刻排尿"避孕}：这些做法对避孕无效，反而可能增加感染风险。
\end{itemize}

\begin{tcolorbox}[colback=pink!5!white,colframe=red!50!black,title=贴心小叮咛]
避孕套的失败，绝大多数不是因为"这东西不靠谱"，而是\textbf{尺寸不合、没有全程使用、或者开了油润滑的玩笑}。记住三句话：\textbf{选对号、从开始就戴、只用水基/硅基润滑}——正确使用时它的避孕成功率可达约 98\%。
\end{tcolorbox}
